% Exact candidate proof for P06.1.
% Not a live-book input before independent and production verification.
% No novelty or priority claim is made.

\begin{proof}[Proof of \Cref{lem:p061-rational-section-tube}]
We prove the first assertion by using the definition of the completed
tensor product.  Fix \(G\in\ker\epsilon\) and \(\eta>0\).  Choose
compatible Banach norms on \(A\), \(B\), and \(P\).  The algebraic tensor
product is dense in \(P\), evaluation by \(\epsilon\) is bounded, and
the supremum seminorm is bounded by any affinoid Banach norm.  We may
therefore choose a finite sum
\[
                       T=\sum_{\ell=1}^N a_\ell\otimes b_\ell
\]
and a positive number \(\delta<\eta\) such that
\[
 \sup_{\mathcal M(P)}|G-T|<\delta,
 \qquad
 \sup_{\mathcal M(B)}|\epsilon(G-T)|<\delta.
\]
Since \(\epsilon(G)=0\), the element
\[
                       b_0=\epsilon(T)=
                       \sum_{\ell=1}^Na_\ell(a)b_\ell
\]
has supremum norm smaller than \(\delta\).  Rewrite
\begin{equation}
\label{eq:p061-tensor-kernel-rewrite}
 T=\sum_{\ell=1}^N
       \bigl(a_\ell-a_\ell(a)\bigr)\otimes b_\ell+1\otimes b_0.
\end{equation}
Every \(a_\ell-a_\ell(a)\) vanishes at \(a\).  Since the valuation of
\(K\) is nontrivial, we can choose arbitrarily small positive radii in
\(|K^\times|\).  The simultaneous rational domains cut out by weak
inequalities at such radii are strict affinoid neighborhoods of \(a\).
Hence there is a common strict affinoid neighborhood \(U\) of \(a\) on
which
\[
 \sup_U|a_\ell-a_\ell(a)|\,
       \sup_{\mathcal M(B)}|b_\ell|<\delta
 \qquad(1\leq\ell\leq N).
\]
For every \(w\in U\times_K\mathcal M(B)\), the ultrametric triangle
inequality and \eqref{eq:p061-tensor-kernel-rewrite} give
\[
 |T(w)|\leq
 \max\!\left\{
   \sup_{\mathcal M(B)}|b_0|,
   \max_{1\leq\ell\leq N}
   \sup_U|a_\ell-a_\ell(a)|\,
   \sup_{\mathcal M(B)}|b_\ell|
 \right\}<\delta.
\]
Applying the same inequality to \(G=(G-T)+T\) gives
\[
                  \sup_{U\times_K\mathcal M(B)}|G|<\delta<\eta.
\]
This proves the first assertion.

For the second assertion, let \(\Omega\) be a neighborhood of \(s_a(z)\)
inside a product affinoid
\(\mathcal M(A)\times_K\mathcal M(B)=\mathcal M(P)\).  Choose a basic
neighborhood
\[
 \Omega_0=
 \bigcap_{i=1}^p\{|F_i|<r_i\}
 \cap
 \bigcap_{j=1}^q\{|H_j|>s_j\}
 \subseteq\Omega
\]
of \(s_a(z)\), where \(F_i,H_j\in P\) and \(r_i,s_j>0\).  For every
upper condition choose
\[
                  |\epsilon(F_i)(z)|<\beta_i<r_i,
\]
and for every lower condition choose
\[
                  s_j<\alpha_j<|\epsilon(H_j)(z)|.
\]
Continuity and goodness give a common strict affinoid neighborhood \(V\)
of \(z\) on which
\[
       |\epsilon(F_i)|<\beta_i
       \quad\text{and}\quad
       |\epsilon(H_j)|>\alpha_j
\]
for all \(i,j\).  For each of the finitely many functions \(L=F_i,H_j\)
write
\begin{equation}
\label{eq:p061-tensor-section-split}
                L=1\otimes\epsilon(L)
                  +\bigl(L-1\otimes\epsilon(L)\bigr).
\end{equation}
The second summand lies in \(\ker\epsilon\).  Apply the first assertion
simultaneously, choosing the error strictly smaller than \(r_i\) for the
summand associated to \(F_i\), and strictly smaller than \(\alpha_j\) for
the summand associated to \(H_j\).  This gives one strict affinoid
neighborhood \(U\) of \(a\).  On \(U\times_KV\), the ultrametric
inequality gives \(|F_i|<r_i\).  For \(H_j\), the error is strictly
smaller than \(|\epsilon(H_j)|\), so
\[
                         |H_j|=|\epsilon(H_j)|>\alpha_j>s_j.
\]
Thus \(U\times_KV\subset\Omega_0\subseteq\Omega\).  This proves the
asserted neighborhood basis without identifying the topology of a
Berkovich product with the ordinary product topology.

For the last assertion, replace \(A\), \(B\), \(P\), and \(\epsilon\) in
the first assertion by
\[
 \mathcal O(U_0),\qquad \mathcal O(V),\qquad
 \mathcal O(U_0)\widehat\otimes_K\mathcal O(V),\qquad
 \epsilon_a\widehat\otimes1,
\]
respectively.  Put \(f_a=s_a^*f\in\mathcal O(V)^\times\) and
\[
                       E=f(1\otimes f_a^{-1})-1.
\]
Then \(E\in\ker\epsilon\).  The first assertion, applied with
\(\eta=1\), lets us shrink \(U_0\) so that \(|E|<1\) on
\(U_0\times_KV\).  Hence \(|1+E|=1\), which is precisely
\eqref{eq:p061-unit-constant-transverse}.
\end{proof}

\begin{proof}[Proof of \Cref{lem:p061-parameterized-rigid-point}]
Choose strict affinoid neighborhoods \(S_1\subset S\) of \(a\) and
\(Z_1\subset Z\) of \(z\).  The neighborhood-basis assertion of
\Cref{lem:p061-rational-section-tube}, applied after restricting to these
affinoids, supplies strict affinoid neighborhoods \(S_0\subset S_1\) and
\(Z_0\subset Z_1\) such that \(S_0\times_KZ_0\) is contained both in the
domain of \(u\) and in the pure-dimensional neighborhood from the
hypothesis.  We may therefore replace the ambient product by this
compact good strict pure-dimensional affinoid product before invoking
the results of Gubler--Jell--Rabinoff.

By \cite[Proposition~5.8]{GJR26}, after replacing this product by a
smaller analytic neighborhood \(\Omega\) of \(s_a(z)\), there are
finitely many \(\lambda_1,\ldots,\lambda_m\in\mathbf R\) and integral PL
functions \(h_1,\ldots,h_m\) on \(\Omega\) such that
\begin{equation}
\label{eq:p061-real-pl-decomposition}
                              u=\sum_{j=1}^m\lambda_jh_j.
\end{equation}
Choose a strict affinoid neighborhood \(W\) of \(s_a(z)\) contained in
\(\Omega\).  It is compact, good, strict, and pure-dimensional, so the
standing hypotheses of Section~8 of \cite{GJR26} hold on \(W\).  Apply
\cite[Proposition~8.3]{GJR26} to each \(h_j|_W\).  After a further
shrinking, there are smooth tropicalization maps
\[
 \tau_j=(\log|f_{j1}|,\ldots,\log|f_{jN_j}|)
\]
defined by analytic units on strict affinoid neighborhoods of
\(s_a(z)\), and piecewise integral-affine functions \(F_j\), such that
\[
                              h_j=F_j\circ\tau_j.
\]
This is the precise passage from a G-local unit presentation to a
presentation on an actual neighborhood; no member of the original
G-cover is being treated as a topological neighborhood.

Apply the neighborhood-basis assertion once more to place a product
\(U\times_KV\subset S_0\times_KZ_0\) inside all the finitely many
neighborhoods just obtained.  List all the \(f_{j\ell}\) as
\(f_1,\ldots,f_N\).  By
\eqref{eq:p061-unit-constant-transverse}, after one common shrinking of
\(U\),
\[
                     |f_k|=p_Z^*|s_a^*f_k|
                     \quad(1\leq k\leq N)
\]
on \(U\times_KV\).  Every \(\tau_j\), and hence every
\(F_j\circ\tau_j\), is therefore pulled back from the section.  Equation
\eqref{eq:p061-real-pl-decomposition} gives
\eqref{eq:p061-parameterized-rigid-point}.  Finally,
\(s_a^*u|_V\) is \(\mathbf R\)-PL by pullback functoriality
\cite[Lemma~5.4(1)]{GJR26}.
\end{proof}

\begin{proof}[Proof of
\Cref{prop:p061-curve-product-normalization-descent}]
The assertion is local on \(C\times_KZ\).  Fix a point \(w\), and write
\(c\in C\) and \(z\in Z\) for its images.  If \(C\) is normal near
\(c\), then \(\nu\) is an isomorphism near \(c\), and there is nothing
to prove.

Assume that \(c\) is nonnormal.  By
\cite[Lemma~3.9(2)]{GR25Classical}, the nonnormal locus is discrete and
consists of rigid points.  Since \(K\) is algebraically closed, \(c\) is
\(K\)-rational and \(w=s_c(z)\).  After shrinking \(C\) around \(c\),
normalization is finite and
\[
                         \nu^{-1}(c)=\{c_1,\ldots,c_r\},
\]
where every \(c_i\in C^\nu(K)\).
The field \(K\) is perfect, so
\cite[Lemma~3.9(1)]{GR25Classical} shows that \(C^\nu\) is quasi-smooth;
as the normalization of a curve, it is pure of dimension one.  If \(Z\)
is pure of dimension \(d\) near \(z\), the strict analytic product
dimension formula shows that \(C^\nu\times_KZ\) is pure of dimension
\(d+1\) near every \((c_i,z)\).  Thus the pure-dimensional hypothesis in
\Cref{lem:p061-parameterized-rigid-point} is satisfied.

Put \(\widetilde g=q^*g\).  Apply
\Cref{lem:p061-parameterized-rigid-point} at every \((c_i,z)\).  After
replacing the finitely many parameter neighborhoods by a common strict
affinoid neighborhood \(V\) of \(z\), there are strict affinoid
neighborhoods \(U_i\) of \(c_i\) and \(\mathbf R\)-PL functions \(r_i\)
on \(V\) such that
\begin{equation}
\label{eq:p061-branchwise-pullback}
                    \widetilde g|_{U_i\times_KV}=p_Z^*r_i.
\end{equation}
Restriction to the section \(\{c_i\}\times_KV\) shows that
\[
                     r_i(v)=g(c,v)\qquad(v\in V).
\]
Hence all \(r_i\) are one and the same \(\mathbf R\)-PL function \(r\).

Choose a strict affinoid neighborhood \(C_0\) of \(c\) contained in the
region under consideration, and put
\[
 T=C_0\times_KV,\qquad \widetilde T=q^{-1}(T).
\]
Inside the compact space \(\widetilde T\), let
\[
 \Omega=\widetilde T\cap
         p_{C^\nu}^{-1}\!\left(\bigcup_{i=1}^r\operatorname{Int}(U_i)\right).
\]
After shrinking the \(U_i\) inside their original expression
neighborhoods if necessary, \(\Omega\) is an open neighborhood of
\[
 |q|^{-1}(\{c\}\times|V|)
       =\coprod_{i=1}^r s_{c_i}(|V|),
\]
and \eqref{eq:p061-branchwise-pullback} gives
\(\widetilde g=p_Z^*r\) on \(\Omega\).

The set \(F=\widetilde T\setminus\Omega\) is compact.  The finite map
\(q:\widetilde T\to T\) is continuous and surjective, so \(q(F)\) is
compact and hence closed in the Hausdorff affinoid space \(T\).  Thus
\[
                              W=T\setminus q(F)
\]
is an open neighborhood of \(\{c\}\times_KV\), and
\(q^{-1}(W)\subset\Omega\).  We have
\[
                  q^*g=q^*p_Z^*r\quad\text{on }q^{-1}(W).
\]
Surjectivity of \(q\) gives \(g=p_Z^*r\) on \(W\).  The right-hand side
is \(\mathbf R\)-PL by \cite[Lemma~5.4(1)]{GJR26}.  This proves the
assertion near \(w\), and hence everywhere.
\end{proof}

\begin{proof}[Proof of \Cref{thm:p061-curve-product-ph-rpl}]
The question is local, so we may shrink \(Y\) around any chosen point
until it is pure-dimensional.  Let \(K'\) be the completion of an
algebraic closure of \(K\), and attach a prime to every object after
scalar extension to \(K'\).  Geometric reducedness makes \(C'\) reduced.
Let
\[
                            \nu':(C')^\nu\longrightarrow C'
\]
be its normalization.  Since \(K'\) is algebraically closed and
therefore perfect, \cite[Lemma~3.9(1)]{GR25Classical} makes
\((C')^\nu\) quasi-smooth.  Consequently
\((C')^\nu\times_{K'}Y'\) is strict and quasi-smooth.  It is good because
\(\nu'\times\mathrm{id}_{Y'}\) is finite over the good space \(X'\):
the inverse image of a strict affinoid neighborhood under a finite map
is strict affinoid.  Since \((C')^\nu\) is pure of dimension one and
\(Y'\) is pure-dimensional after the initial shrinking, their product is
pure-dimensional as well.

The scalar extension \(h'\) is classically pluriharmonic by
\cite[Proposition~7.2(6)]{GR25Classical}.  Its pullback along
\(\nu'\times\mathrm{id}_{Y'}\) is classically pluriharmonic by
\cite[Proposition~7.2(7)]{GR25Classical}.  The production-verified
quasi-smooth regularity theorem
\Cref{thm:p061-quasismooth-ph-rpl} therefore shows that this pullback is
\(\mathbf R\)-PL.  Apply
\Cref{prop:p061-curve-product-normalization-descent} over \(K'\) to
obtain
\[
                              h'\in\operatorname{PL}_{\mathbf R}(X').
\]
Ground-field descent \cite[Corollary~5.12]{GJR26} now gives
\(h\in\operatorname{PL}_{\mathbf R}(X)\).

Both \(h\) and \(-h\) are classically psh.  For an \(\mathbf R\)-PL
function on a good strict space, \cite[Theorem~4.3]{GR25Classical}
identifies classical psh with semipositivity.  Thus both signs are
semipositive, which is harmonicity in
\cite[Definition~7.2]{GJR26}.

We finally justify the asserted open-subspace version.  After extension
to \(K'\), a point whose projection to \(C'\) is normal has a neighborhood
on which normalization is an isomorphism.  If its projection \(c\) is
nonnormal, then \(c\) is \(K'\)-rational.  The neighborhood-basis part of
\Cref{lem:p061-rational-section-tube} places a product affinoid
\(U\times_{K'}V\) containing the point inside the given open analytic
subspace.  The normalization-descent and saturation argument above then
applies on that product neighborhood.  Thus the proof is local in the
analytic topology in precisely the form needed for open analytic
subspaces.  No assertion about an arbitrary analytic domain at one of its
relative-boundary points is made here.
\end{proof}

\begin{proof}[Proof of \Cref{cor:p061-finite-curve-products-ph-rpl}]
Let \(K'\) be the completion of an algebraic closure of \(K\), use a
prime for scalar extension, and let
\(\nu_i:(C_i')^\nu\to C_i'\) be the normalization.  After shrinking \(Y\)
locally, assume that it is pure-dimensional.  For \(0\leq j\leq m\), set
\[
 X_j=
 C_1'\times_{K'}\cdots\times_{K'}C_j'
 \times_{K'}(C_{j+1}')^\nu\times_{K'}\cdots
 \times_{K'}(C_m')^\nu\times_{K'}Y'.
\]
Let \(\rho_j:X_j\to X_m\) normalize the factors with index greater than
\(j\), and put \(h_j=\rho_j^*h'\).  Every \(X_j\) is finite over the good
space \(X_m\), hence good; it is also strict.  The space \(X_0\) is
quasi-smooth, so functoriality of classical pluriharmonicity and
\Cref{thm:p061-quasismooth-ph-rpl} give
\[
                         h_0\in\operatorname{PL}_{\mathbf R}(X_0).
\]

For \(1\leq j\leq m\), let \(q_j:X_{j-1}\to X_j\) be induced by
\(\nu_j\) on the \(j\)-th factor and by the identity on every other
factor.  Then
\[
                             q_j^*h_j=h_{j-1}.
\]
The parameter product formed by the other curve factors and \(Y'\) is
good and strict locally: at a given point, choose strict affinoid
neighborhoods in all of its factors and take their affinoid product.  It
is pure-dimensional there, since every reduced analytic curve is pure of
dimension one, every normalized curve is pure of dimension one, and
\(Y'\) is locally pure-dimensional.  Thus
\Cref{prop:p061-curve-product-normalization-descent}, with \(C=C_j'\),
applies at each point.  Induction on \(j\) gives
\[
                         h_j\in\operatorname{PL}_{\mathbf R}(X_j)
                         \qquad(0\leq j\leq m).
\]
In particular, \(h_m=h'\) is \(\mathbf R\)-PL.  Ground-field descent
\cite[Corollary~5.12]{GJR26} gives the assertion over \(K\).
Harmonicity follows exactly as in the proof of
\Cref{thm:p061-curve-product-ph-rpl}.

It remains to justify the open-subspace assertion.  It is enough to work
after extension to \(K'\), prove \(\mathbf R\)-piecewise linearity at
every point upstairs, and then use
\cite[Corollary~5.12]{GJR26}.  Fix a point of the given open subspace and
let \(I\subseteq\{1,\ldots,m\}\) be the indices for which its projection
\(c_i\in C_i'\) is nonnormal.  Every \(c_i\) with \(i\in I\) is
\(K'\)-rational.  If \(I\) is empty, the open subspace is already locally
quasi-smooth and \Cref{thm:p061-quasismooth-ph-rpl} applies directly.
Assume henceforth that \(I\ne\varnothing\).  For \(i\notin I\), first
shrink the corresponding curve coordinate inside its normal locus.
Starting with the resulting open neighborhood and applying
the neighborhood-basis part of
\Cref{lem:p061-rational-section-tube} successively to these rational
coordinates, we obtain strict affinoid neighborhoods \(U_i\ni c_i\) for
\(i\in I\) and a strict affinoid neighborhood \(V\) in the product of
the remaining coordinates such that
\[
                       \prod_{i\in I}U_i\times_{K'}V
\]
is contained in the given open subspace.  By
\cite[Lemma~3.9(1)]{GR25Classical}, the chosen curve neighborhoods for
\(i\notin I\), which lie in the normal locus and identify with open
subspaces of the normalization, are quasi-smooth.  Their product with
\(Y'\) is quasi-smooth, and therefore so is its analytic domain \(V\).
The normalization-and-induction argument above, now with the reduced
possibly singular curve domains \(U_i\) for \(i\in I\) and the
quasi-smooth factor \(V\), proves \(\mathbf R\)-piecewise linearity on a
neighborhood of the chosen point.  This proves the assertion at every
point upstairs and completes the descent.
\end{proof}

\begin{proof}[Verification of \Cref{ex:p061-nodal-cuspidal-products}]
Both displayed affinoid curves are geometrically reduced in
characteristic zero.  Their only singular point is respectively the node
or the cusp.  The assertion follows from
\Cref{thm:p061-curve-product-ph-rpl,cor:p061-finite-curve-products-ph-rpl}.
\end{proof}
